\section{神经隐式表示（Neural Implicit Representations）}

\subsection{从传统到神经网络}

2019 年前后，研究者意识到一个关键洞察：深度网络本质上就是在学习一个复杂的\textbf{隐式函数}——输入空间大（像素），输出空间小（类别）。那么，为何不直接用神经网络来\textbf{表示}一个 3D 物体的隐式函数？

\begin{figure}[H]
\centering
\includegraphics[width=0.95\textwidth]{slides-images/slide-035.jpg}
\caption{深度隐式函数（Deep Implicit Functions）：神经网络学习将 $(x,y,z)$ 映射到占用概率或距离值\protect\footnotemark}
\end{figure}
\footnotetext{Slides 第35页。}

\begin{importantbox}{2019年的四篇里程碑论文}
2019年几乎同时出现了四篇提出相同思想的论文：
\begin{itemize}
    \item \textbf{DeepSDF}（Park et al.）：学习有符号距离函数
    \item \textbf{Occupancy Networks}（Mescheder et al.）：学习占用概率
    \item \textbf{IM-NET}（Chen \& Zhang）：隐式场解码器
    \item \textbf{DISN}（Xu et al.）：深度隐式表面网络
\end{itemize}
核心思想一致：用 MLP 接收 $(x,y,z)$ 坐标作为输入，输出该点是否在物体内部的概率或到表面的距离。
\end{importantbox}

网络结构非常简单：
$$\text{MLP}: (x, y, z) \rightarrow \text{occupancy} \in [0, 1] \quad \text{或} \quad \text{SDF value} \in \mathbb{R}$$

\subsection{从几何到外观：NeRF}

如果隐式函数不仅输出几何（密度），还输出\textbf{颜色}，就可以实现从任意视角的渲染。这就是 2020 年 \textbf{NeRF}（Neural Radiance Fields）的核心思想。

\begin{figure}[H]
\centering
\includegraphics[width=0.95\textwidth]{slides-images/slide-039.jpg}
\caption{NeRF 管线：沿射线采样 3D 点，每个点查询 MLP 得到密度和颜色，通过体渲染合成像素颜色\protect\footnotemark}
\end{figure}
\footnotetext{Slides 第39页。}

$$\text{NeRF MLP}: (x, y, z, \theta, \phi) \rightarrow (\sigma, \mathbf{c})$$

\begin{itemize}
    \item $(x, y, z)$：3D 位置
    \item $(\theta, \phi)$：观察方向
    \item $\sigma$：体密度（density）
    \item $\mathbf{c} = (r, g, b)$：颜色
\end{itemize}

\textbf{体渲染公式}：对于一条射线 $\mathbf{r}(t) = \mathbf{o} + t\mathbf{d}$：

$$\hat{C}(\mathbf{r}) = \sum_{i=1}^{N} T_i \cdot \alpha_i \cdot \mathbf{c}_i, \quad T_i = \prod_{j=1}^{i-1}(1 - \alpha_j), \quad \alpha_i = 1 - e^{-\sigma_i \delta_i}$$

\begin{itemize}
    \item $T_i$：透射率（transmittance），表示光线到达第 $i$ 个采样点前未被遮挡的概率
    \item $\alpha_i$：不透明度（opacity）
    \item $\delta_i$：相邻采样点间距
\end{itemize}

\begin{figure}[H]
\centering
\includegraphics[width=0.95\textwidth]{slides-images/slide-040.jpg}
\caption{NeRF 体渲染过程：沿每条射线采样、查询 MLP、加权求和得到像素颜色\protect\footnotemark}
\end{figure}
\footnotetext{Slides 第40页。}

\begin{importantbox}{NeRF 的训练与推理}
\textbf{训练}：仅需多张已知相机位姿的照片。对每张照片的每个像素发射射线，通过体渲染合成颜色，与真实像素颜色计算 MSE 损失，反向传播更新 MLP 权重。\\
\textbf{推理}：给定新的相机位姿，即可渲染出该视角的图像——实现\textbf{新视角合成}（Novel View Synthesis）。
\end{importantbox}

\begin{figure}[H]
\centering
\includegraphics[width=0.95\textwidth]{slides-images/slide-041.jpg}
\caption{NeRF 的新视角合成效果：从一组输入图像学习 3D 场景，渲染任意新视角\protect\footnotemark}
\end{figure}
\footnotetext{Slides 第41页。}

\subsection{3D Gaussian Splatting}

NeRF 虽然效果出色，但渲染速度慢——需要沿每条射线密集采样大量 3D 点并查询 MLP。很多采样点对应的是空白区域，造成计算浪费。

\textbf{3D Gaussian Splatting}（3DGS，2023）提出了一种加速方案：用一组\textbf{3D 高斯椭球}（Gaussian Blobs）来显式表示场景。

\begin{figure}[H]
\centering
\includegraphics[width=0.95\textwidth]{slides-images/slide-043.jpg}
\caption{3D Gaussian Splatting：用 3D 高斯椭球显式表示场景，只在有物体的区域采样，大幅提速\protect\footnotemark}
\end{figure}
\footnotetext{Slides 第43页。}

\begin{importantbox}{NeRF vs 3D Gaussian Splatting}
\begin{itemize}
    \item \textbf{NeRF}：纯隐式（MLP），密集采样，渲染慢但表示紧凑
    \item \textbf{3DGS}：半显式（高斯点云），稀疏采样，渲染快（可实时）但存储量大
\end{itemize}
3DGS 仍保留隐式查询的优势（每个高斯有密度和颜色属性），但因为知道``哪里有东西''，避免了对空白区域的无效采样。
\end{importantbox}

\begin{figure}[H]
\centering
\includegraphics[width=0.95\textwidth]{slides-images/slide-044.jpg}
\caption{3DGS 重建效果与速度对比\protect\footnotemark}
\end{figure}
\footnotetext{Slides 第44页。}

\subsection{本章小结}

\begin{itemize}
    \item 神经隐式表示用 MLP 作为 3D 物体的隐式函数，输入坐标、输出占用/距离/颜色
    \item NeRF 将隐式表示扩展到辐射场，实现了从 2D 图像学习 3D 场景并渲染新视角
    \item 3D Gaussian Splatting 结合显式点云和隐式属性，实现了接近实时的渲染速度
    \item 这一系列方法打通了 2D 图像数据与 3D 场景理解之间的桥梁
\end{itemize}

%% ========== 3D 生成 ==========

